% Generated by code/interviews-1/03-speed-drills/python/gen_drill.py. Do not edit. French copy by draft/pilot-quant-fr/ivgen_fr.py.
\ivfamily{Fractions en pourcentages}{2, 3, 13, 20, 21, 24, 39, 40, 41, 47}{Connaître les huitièmes et les seizièmes. Pour un dénominateur de 20, 25 ou 50, ramener à 100~: $\tfrac{7}{25} = \tfrac{28}{100}$. Pour 40, prendre la moitié de la valeur sur 20.}
\ivfamily{Décimaux en fractions}{26, 27, 28, 30, 33, 36, 37, 45}{Lire le décimal sur une puissance de 10, puis simplifier~: $0.375$ vaut $\tfrac{375}{1000}$, soit $\tfrac38$. La plupart sont des multiples de $\tfrac18$, $\tfrac1{20}$, $\tfrac1{25}$ ou $\tfrac1{40}$.}
\ivfamily{Variation en pourcentage}{1, 4, 6, 8, 29, 32, 34, 43, 51, 53}{Ajouter ou soustraire le pourcentage de la base, ou multiplier directement~: $15\%$ en hausse s'écrit $\times 1.15$, et $20\%$ en baisse s'écrit $\times 0.8$.}
\ivfamily{En pourcentage de}{11, 15, 19, 22, 23, 25, 35, 52}{Simplifier d'abord, puis ramener à 100~: $36$ sur $120$ vaut $\tfrac{3}{10}$, soit $30\%$.}
\ivfamily{Pourcentages de n'importe quel nombre}{5, 16, 17, 18, 31, 38, 46, 49}{Prendre 1~\% ou 10~\% et monter, ou permuter les deux~: $p\%$ de $n$ égale $n\%$ de $p$.}
\ivfamily{Fractions}{7, 9, 10, 12, 14, 42, 44, 48, 50, 54}{Utiliser le plus petit dénominateur commun. Pour diviser, multiplier par la fraction inverse, et simplifier d'abord.}
